\documentclass[twocolumn]{iopjournal}
\usepackage{amsmath,amssymb,amsfonts,amsthm}
\fancyhfoffset[L]{0mm}
\setlength{\marginparwidth}{0pt}
\setlength{\marginparsep}{0pt}
\renewcommand{\articletype}[1]{%
  {\vspace*{-8mm}\noindent \Large \sf Preprint}\par
  \vspace*{6mm}\noindent{\scriptsize \sf{\bfseries \MakeUppercase{#1}}}}
\fancyhead[L]{}
\fancyhead[R]{}
\hypersetup{pdftitle={The Locked Bimetric Potential and the Constitutive Law}, pdfauthor={Stephen Nelson}, pdfsubject={Gravity Field Dynamics (GFD)}, pdfkeywords={Hassan-Rosen, bimetric gravity, Vainshtein, Galileon, exchange symmetry, constitutive law, MOND, proof}}

\begin{document}
\articletype{Supplementary Material}
\title{The Locked Bimetric Potential and the Constitutive Law}
\author{Stephen Nelson}
\affil{Field Dynamics Research}
\ead{stephen@fielddynamics.org}

\begin{abstract}
Papers~VIII, IX and XII fix the Hassan--Rosen coefficients to
$\beta = (1,0,16,0,1)$ and state that the locked bimetric theory and
the SST are two projections of one geometry, so that the bimetric
theory's quasi-static limit should give $\mu(x) = x/(x+1)$. We show it
does not: Hassan--Rosen's two metrics are not GFD's two chambers.
\emph{Negative:} in the standard
reduction (decoupling limit, static spherical weak field) the response
depends on the enclosed mass only through $M/r^3$, for every $\beta$;
no function of $M/r^3$ equals a function of $x \propto M/r^2$ for all
masses and radii; on the branch that matches the linear solution the
enhancement is bounded, $g_{\rm tot}/g_N \le 1 + 1/\sqrt2 < 1.71$, and
the other real roots give either $g_{\rm tot}/g_N \le 1/3$ or a
mass-independent force growing like $r$; and the potential vanishes at
GFD's parity configuration. Letter~2's Lemma explains why: TetA and
TetB are two roots of one metric, so a potential built from
$\sqrt{g^{-1}f}$ takes the same value on either root
($g^{-1}f = I$) and cannot see the partition. The obstruction is
structural, not specific to the lock.
\emph{Diagnostic:} on the frame deformed along a single leg by a
ratio $w$, every potential is $V = A + Bw$, and the deformed leg's share
$Bw/(A + Bw)$ obeys exchange closure $s(1/w) = 1 - s(w)$ if and only
if $\beta_0 + 2\beta_1 = \beta_4 + 2\beta_3$. It is then exactly
$w/(1+w)$. Palindromic coefficients (the $g\leftrightarrow f$
symmetry) satisfy the condition, the locked set gives $V = 49(1+w)$,
and $\beta_2 = 16$ cancels. With $w = |a|/a_0$, $a$ the acceleration
of the tetrad's time leg, this is $\mu(x)$. The deformation is not a
Hassan--Rosen reference metric, the two shares are parts of one
deformed frame rather than the two roots (on the roots the potential is
98 and 0), and the energy $F$ is not supplied by the potential. So this
reproduces the algebra of the partition, not its two-chamber mechanism.
The constitutive law's origin remains the two-root partition under
exchange (Letter~2).
\end{abstract}

\keywords{Hassan--Rosen, bimetric gravity, Vainshtein, Galileon, exchange symmetry, constitutive law, MOND, proof}

\section*{Setting}

The Hassan--Rosen (HR) interaction is
$m^2\sqrt{-g}\,V(S)$, $V(S) = \sum_{n=0}^4\beta_n e_n(S)$, with
$S = \sqrt{g^{-1}f}$ and $e_n$ the elementary symmetric polynomials of
its eigenvalues~\cite{HR2012}. The lock is
$\beta = (1,0,16,0,1)$~\cite{NelsonVIII,NelsonXII}. The SST's scalar
sector is $-(a_0^2/8\pi G)F(y)$, $F(y) = y - 2\sqrt y + 2\ln(1+\sqrt y)$,
$\mu(x) = F'(x^2) = x/(x+1)$~\cite{NelsonL3,NelsonPqs}. The question,
listed as open in Papers~VIII and IX, is whether the first reduces to
the second. Each step is checked by the accompanying script, 48 checks,
exact algebra~\cite{NelsonScript}, including the decoupling-limit
derivation itself, done independently of~\cite{NelsonBimetric}.

\section*{Part A: the standard reduction}

\textbf{Theorem 1 (the response depends on density).} \emph{In the
decoupling limit, with Stückelberg field $\pi$ and
$\Lambda_3^3 = m^2M_{\rm Pl}$, the $h$--$\pi$ mixing has coefficients
$d_0 = \beta_0+3\beta_1+3\beta_2+\beta_3$,
$d_1 = \beta_1+2\beta_2+\beta_3$, $d_2 = \beta_2+\beta_3$,
$d_3 = \beta_3$; for the lock, $(49, 32, 16, 0)$. Around a static mass,
with $d_3 = 0$, $\lambda = \pi'/r = -sq/(2Ma_1)$ and
$Q = a_2q/(Ma_1^2)$, the $\pi$-equation is
\begin{equation}
Q^2s^3 + 6Qs^2 + 6s - 2 = 0,\qquad g_{\rm tot}/g_N = 1 + s,
\label{eq:screen}
\end{equation}
with $q = M_{\rm enc}/4\pi r^3$. The radius enters only through $q$.}

\begin{proof}
The identity $\sum_n(-1)^n\beta_nY_n(1-\Pi) = \sum_kd_kY_k(\Pi)$, with
$Y_n(X) = \sum_k(-1)^ke_k(X)X^{n-k}$, is checked for general $\beta$
on an exact random symmetric $\Pi$ ($Y_4 = 0$ by Cayley--Hamilton). From
$-\tfrac12h\mathcal Eh + h\cdot\sum_ka_kY_k(\Pi) + h_{00}\rho/2M$ with
$h = \mathrm{diag}(A,B,B,B)(r)$ and $\pi(r)$, all three Euler--Lagrange
equations are $r^{-2}$ times the derivative of a current (checked), so
they integrate exactly. Regularity at the origin gives the
$\pi$-equation
$\tfrac32a_1^2\lambda - (3a_1a_2 + a_3q/4M)\lambda^2 +
(a_2^2 + \tfrac23a_1a_3)\lambda^3 - \tfrac16a_3^2\lambda^5 + a_1q/4M = 0$,
which contains no $r$, and
$g_{\rm tot}/g_N = 1 - 2Ma_1\lambda/q + \tfrac23Ma_3\lambda^3/q$
(checked). Setting $a_3 = 0$ and substituting gives~(\ref{eq:screen}).
\end{proof}

$d_1 = 32$ is the Fierz--Pauli combination and $d_3 = 0$ removes the
quintic mixing; $d_2 = 16$ keeps the cubic Galileon term (Paper~XII
states that both vanish).

\textbf{Theorem 2 (the branches).} \emph{Every real root
of~(\ref{eq:screen}) is $s = 2/D(v)$, $Q = vD(v)/2$, with
$D = v^2 + 6v + 6$.
\begin{enumerate}
\item The branch that matches the linear solution ($v\to0$: $s = 1/3$,
$g_{\rm tot}/g_N = 4/3$, vDVZ) is screened for $Q>0$ ($s\to0$). For
$Q<0$ it turns at $Q_* = 2 - 2\sqrt2$, where $s = 1/\sqrt2$, and does
not continue past it. On it, $g_{\rm tot}/g_N \le 1 + 1/\sqrt2 < 1.71$.
\item Where $D<0$, $s\le -2/3$, so $g_{\rm tot}/g_N\le 1/3$.
\item As $D\to0^+$ ($v\to-3\pm\sqrt3$), $Q\to0$ and $s\simeq v/Q\to\infty$.
There $g_{\rm tot}\simeq g_Nv/Q \propto (M/r^2)(r^3/M) = r$,
independent of the mass.
\end{enumerate}}

\begin{proof}
Substituting the parametrisation solves~(\ref{eq:screen}) identically;
the turning points are $v = -2\pm\sqrt2$, $\min D = -3$, and the
roots of $D$ are $-3\pm\sqrt3$, where $Q = 0$ (all checked).
\end{proof}

The unbounded roots of~(\ref{eq:screen}) do not match the linear solution at large
$r$, and their force grows like $r$ whatever the mass: a
self-accelerating-type branch, not the deep-MOND form
$\sqrt{GMa_0}/r$. So no branch gives flat rotation curves. Theorem~3
shows this for every branch at once.

\textbf{Theorem 3 (no change of variable).} \emph{No function of
$M/r^3$ equals a function of $x = GM/(a_0r^2)$ for all $M$ and $r$,
unless both are constant. Since every root of the $\pi$-equation is a
function of $q$ alone (Theorem~1), the standard response cannot be
$1/\mu(x)$ on any branch, for any $\beta$.}

\begin{proof}
Along $M = Lr^3$ the density is fixed while $x = GLr/a_0$ varies, and
$\mathrm d[1/\mu(x)]/\mathrm dr \neq 0$ (checked). Equivalently, the
Vainshtein radius scales as $M^{1/3}$ and the MOND radius
$\sqrt{GM/a_0}$ as $M^{1/2}$.
\end{proof}

With $m\sim H_0$, galaxies also lie well
inside the Vainshtein radius ($\sim$260~kpc at $10^{10}M_\odot$), where
$\mu_{\rm eff}\to1$~\cite{NelsonBimetric}. Dimensionally, the Galileon
variable is $\Phi/(mr)^2$ and the MOND variable $\Phi/(mr)$: both small
quantities $\Phi/c^2$ and $rH_0/c$ are $\sim10^{-6}$ at galactic radii,
which is why $a_0\sim cH_0$ is natural, but the potential's strongest
interactions select the wrong power.

\textbf{Theorem 4 (the parity configuration and the polynomial).}
\emph{At $S = P = \mathrm{diag}(1,-1,-1,-1)$, $e_n(P) = (1,-2,0,2,-1)$
and $V(P) = \beta_0 - \beta_4 = 0$. On $S = \mathrm{diag}(1,s,s,s)$ the
locked potential is $1 + 48s + 48s^2 + s^3$ (98, 1, 0 at
$s = +1, 0, -1$), not the coupling polynomial $1 + 4s + 16s^2$
(21, 17, 13).}

\begin{proof}
Direct evaluation (checked).
\end{proof}

So GFD's two roots (one metric, $S = P$) and HR's two metrics are
different objects, and the potential's values are not the polynomial's.

\textbf{Theorem 5 (Hassan--Rosen is root-blind).} \emph{TetA $= e$
and TetB $= Pe$ give the same metric, $g = e^{\mathsf T}\eta e =
(Pe)^{\mathsf T}\eta(Pe)$, while $\det e$ changes sign. Hence, with $f$
built from TetB, $g^{-1}f = I$, and the HR potential takes the same
value on either root. Only the vielbein form $e^{-1}(Pe) = P$ sees the
second root, and there $V(P) = 0$.}

\begin{proof}
$P^{\mathsf T}\eta P = \eta$; checked for a general static tetrad.
\end{proof}

The two roots are the two square-root branches of $g^{-1}f = I$: $P$
satisfies $P^2 = I$ but is not the principal root (its spatial
eigenvalues are $-1$), and $e^{\mathsf T}\eta(Pe)$ is symmetric, so the
vielbein form with $S = P$ meets the symmetric-vielbein condition
(checked). The metric form, defined with the principal root, sees only
$S = I$.

This is Letter~2's Lemma~\cite{NelsonL2}: objects built from the
metric (the metric, a torsion-free connection) cannot tell TetA from
TetB, and only the torsion sector can. The HR interaction is built from
the metrics and contains no torsion, so it cannot carry the TetA/TetB
partition. That is the reason for Theorems~1--4.

\section*{Part B: a deformed frame (diagnostic)}

\textbf{Theorem 6 (exchange closure on a deformed frame).} \emph{On
$S = \mathrm{diag}(1,1,1,w)$, for every $\beta$,
\[
\begin{aligned}
V &= A + Bw,\\
A &= \beta_0+3\beta_1+3\beta_2+\beta_3,\\
B &= \beta_1+3\beta_2+3\beta_3+\beta_4 .
\end{aligned}
\]
The deformed leg's share $s(w) = Bw/(A + Bw)$ satisfies
$s(1/w) = 1 - s(w)$ if and only if $A = B$, that is
\begin{equation}
\beta_0 + 2\beta_1 = \beta_4 + 2\beta_3 ,
\label{eq:cond}
\end{equation}
and then $s(w) = w/(1+w)$ exactly.}

\begin{proof}
$s(1/w) - [1 - s(w)] = -w(A-B)(A+B)/[(A+Bw)(Aw+B)]$ (checked), which
vanishes iff $A = B$ since $A + B > 0$.
\end{proof}

The share is not an arbitrary choice: $Bw/(A+Bw) = w\,\mathrm d\ln
V/\mathrm dw$, the elasticity of $V$ on this family (checked).

\textbf{Corollary 1 (the lock and $\beta_2$).} \emph{HR's $g\leftrightarrow f$ symmetry
$\det S\,V(S^{-1}) = V(S)$ holds iff $\beta$ is palindromic (the
difference is $\sum_n(\beta_{4-n}-\beta_n)e_n(S)$). Palindromic $\beta$
satisfies (\ref{eq:cond}), and then
$V = (\beta_0 + 4\beta_1 + 3\beta_2)(1+w)$. The locked set gives
$V = 49(1+w)$, and $\beta_2$ cancels from $s$. On this family
$V/V(I) = \det[(I+S)/2]$, the volume of the midpoint frame.}

\begin{proof}
All identities checked exactly; $(1,0,16,0,2)$ shows that
(\ref{eq:cond}) is a real condition.
\end{proof}

The share's form follows from the $g\leftrightarrow f$ symmetry,
not from the value 16.

\textbf{Remark (not the two chambers).} The two shares $w/(1+w)$ and
$1/(1+w)$ sum to 1, have odds $w$, and swap under $w\to1/w$: the
algebra of the partition. But they are the deformed leg and the three
undeformed legs of one frame, not TetA and TetB. On the actual roots
the locked potential is $V(I) = 98$ and $V(P) = 0$, and on the path
between them, $\mathrm{diag}(1,s,s,s)$, it is $(1+s)(s^2+47s+1)$, with
no exchange symmetry between the ends (checked). So this is not the
two-chamber mechanism of~\cite{NelsonExch,NelsonL2}.

\textbf{Theorem 7 (the time leg, and what the potential does not
supply).} \emph{(i) With $w = |a|/a_0$, $a$ the acceleration of the
tetrad's time leg, a static field has $c^2|a| = |\nabla\Phi|$
exactly~\cite{NelsonTL}, so $w = x$ and $s = \mu(x)$. (ii) The
auxiliary energy
\[
E(u;y) = (1 - e^{-u})\,y - 2(\sinh u - u),
\]
with $u = \ln(V/V_0)$,
is stationary at $V/V_0 = 1 + \sqrt y$, where its value is $F(y)$ and
its slope $\partial E/\partial y$ is $\mu$. (iii) Its $u$-dependence
contains no $\beta$.}

\begin{proof}
All checked exactly.
\end{proof}

So the potential supplies the variable $V/V_0 = 1 + x$; the energy
$F$, and with it the dynamics, comes from the SST's own scalar term.

\section*{Conclusion and scope}

\begin{itemize}
\item The locked HR theory does not reduce to $\mu(x)$ by the standard
reduction (Theorems~1--3), and no choice of $\beta$ would change that.
Papers~VIII, IX and XII overstate the equivalence.
\item What the lock does fix is bimetric physics: the de~Sitter
background $\Lambda = m^2(1+48c^2)$, the Fierz--Pauli mass, the Higuchi
bound and the screening cubic~(\ref{eq:screen})~\cite{NelsonBimetric}.
\item HR is root-blind (Theorem~5), so its two metrics are not GFD's
two chambers.
\item On a frame deformed along the time leg's acceleration the
potential reproduces the partition's algebra (Theorem~6, Corollary~1,
Theorem~7), but not its mechanism: the shares are not the two roots,
the frame is not an HR reference metric, and $F$ is not in $V$. The
constitutive law's origin remains the two-root partition under
exchange~\cite{NelsonL2,NelsonExch}.
\item Letter~2 locates the partition in the torsion sector. A
dual-chamber bridge would therefore couple connections, not metrics;
that is the subject of~\cite{NelsonTwoConn}.
\item Not done here: the full bimetric (dynamical $f$) decoupling
limit, which keeps the linear enhancement below $4/3$; and a
derivation of the deformation $w = |a|/a_0$ from an action.
\end{itemize}

\begin{thebibliography}{99}
\bibitem{HR2012} S.~F.~Hassan and R.~A.~Rosen, Bimetric gravity from
ghost-free massive gravity, \textit{JHEP} \textbf{02} (2012) 126.
\bibitem{NelsonVIII} S.~Nelson, \textit{GFD and Bimetric Gravity},
Paper~VIII (2026).
\bibitem{NelsonXII} S.~Nelson, \textit{Dual-tetrad covariance and the
Hassan--Rosen lock}, Paper~XII (2026).
\bibitem{NelsonL2} S.~Nelson, \textit{The Other Half of Gravity},
Letter~2 (2026), CQG-115447.
\bibitem{NelsonTwoConn} S.~Nelson, \textit{Two connections and the
dual chamber} (2026), research note,
\texttt{research/two-connection-\linebreak dual-chamber/}.
\bibitem{NelsonL3} S.~Nelson, \textit{Field Closure of the Tetrad
Frame}, Letter~3 (2026).
\bibitem{NelsonPqs} S.~Nelson, \textit{The SST Covariant Completion in
Its Quasi-Static Limit} (2026), \texttt{proofs/sst-quasi-static-limit/}.
\bibitem{NelsonTL} S.~Nelson, \textit{The SST Covariant Completion: $y$
From the Acceleration of the Time Leg} (2026),
\texttt{proofs/\linebreak sst-covariant-time-leg/}.
\bibitem{NelsonExch} S.~Nelson, \textit{The constitutive law from
exchange symmetry} (2026),
\texttt{proofs/constitutive-law-\linebreak exchange-symmetry/}.
\bibitem{NelsonBimetric} S.~Nelson, \textit{Does the locked bimetric
potential project onto $\mu(x)$?} (2026),
\texttt{research/bimetric-projection/}.
\bibitem{NelsonScript} S.~Nelson,
\texttt{verify\_bimetric\_lock\_\linebreak constitutive\_law.py} (2026),
\texttt{proofs/bimetric-lock-\linebreak constitutive-law/}.
\end{thebibliography}
\end{document}
